% s9: what the theory contains. ASCII only.
\section{What the theory contains}

Jaffe and Witten ask for local fields for the gauge-invariant polynomials in the curvature, agreement with perturbation
theory at short distance, a theory that is not free, and all of it for every compact simple group~\cite{JaffeWitten}.

\textbf{Local fields.} Traces of products of $F$ and its covariant derivatives are gauge invariant; graded by dimension
($F$ counts 2, each derivative 1), one of the first is $\mathrm{tr}\,F^2$, the action-density scalar, not the energy density $T_{00}$.
A \emph{renormalized local field} is the limit of a lattice transcription (for
$\mathrm{tr}\,F^2$, the six plaquettes at a point) after a finite mixing with fields of equal or lower dimension; its correlations must be tempered distributions
(continuous on Schwartz test functions) obeying the additional growth condition OS0$'$ that reconstruction
needs~\cite{OS75}. Perturbative renormalization is established~\cite{Reisz88,Reisz89}. At positive flow time, flowed
gauge-field observables need no \emph{additional} renormalization after the underlying theory is renormalized.
The local stress tensor requires matching and a small-flow-time limit~\cite{Luscher10,LuscherWeisz11,Suzuki13}. The nearest
complete non-perturbative analogue is the two-dimensional Gross--Neveu model~\cite{GawedzkiKupiainen85,FMRS86}; nothing
comparable is proved in four dimensions.

\begin{claim}{CONJECTURE}{H6: local fields, and agreement with perturbation theory}\label{b:h6}
On the sphere family and its $\mathbb R^4$ end, every gauge-invariant polynomial in $F$ and its covariant derivatives,
$\mathrm{tr}\,F^2$ first, is a renormalized local field (tempered, OS0$'$), with finite mixing and normalization fixed
by a declared prescription and Ward identities, that reconstructs
to a local covariant operator-valued distribution. A conserved stress tensor generates translations, with trace
$\Theta=-\frac{\beta(g)}{2g}[O]$, $\beta(g)=-\frac{b_0}{16\pi^2}g^3+O(g^5)$. Here the Euclidean trace convention
matches the positive running-plaquette coefficient below, and $O=\sum_{a,\mu,\nu}\mathcal F^a_{\mu\nu}\mathcal F^a_{\mu\nu}$
uses canonically normalized fields with Lagrangian $O/4$. This fixes the sign and normalization instead of silently
identifying $O$ with a differently normalized matrix trace. Vacuum subtraction is understood. At short distance every
correlation function agrees with renormalized perturbation theory in the running coupling, asymptotically to all orders.
\evidence{The perturbative ingredients are established~\cite{GrossWilczek73,Politzer73,ACD77,CDJ77,Luscher10,LuscherWeisz11,Suzuki13};
my sphere coefficient $b_0/8\pi^2$ (Statement~\ref{b:b0}) is the logarithm whose cutoff derivative is the one-loop trace
anomaly~\cite{ClosedSet}.}
\lines{25, 14}
\end{claim}

On the lattice there is an exact starting point for the trace anomaly: varying the inverse coupling in the Wilson measure
differentiates every expectation by minus its correlation with the action, and along the line of constant physics that
derivative is the response to a change of scale.

\begin{claim}{PROVED}{The running-plaquette insertion (lattice identity)}\label{b:anomaly}
On every finite lattice, for parameter-independent $X$, $\partial_{(4/g_0^2)}\langle X\rangle=-\langle X;S_W\rangle$,
$S_W=\sum_p(1-\tfrac12\mathrm{tr}\,U_p)$.
Along the line of constant physics this makes the anomaly density $\Theta_a(x)=-\frac{d(4/g_0^2)}{d\ln a}a^{-4}
\sum_{p\text{ based at }x}(s_p-\langle s_p\rangle)$, with six positively oriented plaquettes per site and
$-\frac{d(4/g_0^2)}{d\ln a}=\frac{b_0}{2\pi^2}+O(g_0^2)$ at one loop. Since
$\sum_{p\text{ based at }x}s_p=a^4g_0^2O/16+\cdots$, the coefficient agrees with $-\beta(g)O/(2g)$.
The exact identity fixes an integrated insertion, not all local contact terms or matching. Its continuum
identification as the stress-tensor trace belongs to H6.
\evidence{Differentiation under the integral; factor and plaquette count verified by the referee pass~\cite{ClosedSet}.}
\lines{1, 25}
\end{claim}

\textbf{Non-triviality.} A free massive theory also has a gap. My first criterion (a nonzero connected four-point function of
$\mathrm{tr}\,F^2$) holds in free theories too, and I withdrew it~\cite{ClosedSet}. The fingerprint I now seek is logarithmic
running: by H6, $|x|^8\langle\Theta(x)\Theta(0)\rangle\sim C_O(N^2-1)/\bigl(16\ln^2(1/|x|\Lambda)\bigr)$ as
$|x|\to0$, an inverse square logarithm, unlike a nonzero dimension-four Wick polynomial in finitely many canonical
free fields with a nonvanishing leading scaling term. A \emph{generalized free field}, Gaussian with an arbitrary
two-point function, could still mimic that two-point fingerprint.

I expected the leading three-point function to distinguish it. That particular prediction is now settled against
me: the free Wick triangle is $T_O=0$ at distinct points in Euclidean four-space. The covariance exchanges the
self-dual and anti-self-dual halves of a two-form. Three exchanges cannot close within either half, so their
trace vanishes. This is an exact geometric cancellation, not a numerical near-zero. It says nothing about
contact terms or the higher-order interacting three-point function.

Four exchanges can return to their starting half. The four-point alternative already in my note is nonzero:
for $n=\dim\mathfrak g$ free colours and $O=\sum_{a,\mu,\nu}:F^a_{\mu\nu}F^a_{\mu\nu}:$, at separation $L>0$,
\[
\langle O(0)O(Le_0)O(2Le_0)O(3Le_0)\rangle_c
=\frac{676n}{27\pi^8L^{16}}>0.
\]
The subscript means connected: the disconnected pairings have been removed. With H6's matching this gives a
nonzero leading connected four-point function of $\Theta$ at weak running coupling, proportional to
$(b_0\bar g^2/32\pi^2)^4$. It is a test of whether \emph{this field} is Gaussian, not a restoration of my withdrawn
claim about the whole theory.

\begin{claim}{CONDITIONAL}{Two tests of non-triviality}\label{b:nontrivial}
If H6 holds, the logarithmic two-point fingerprint excludes the stated finite canonical Wick-polynomial comparison,
and the nonzero connected four-point function excludes $\Theta$ itself being a generalized free field.
They do not alone classify the full observable algebra: a non-Gaussian composite of Gaussian fields can have
a nonzero connected four-point function. Non-triviality of the reconstructed theory remains an H6 identification.
\evidence{The $1/16$ coefficient follows from one-loop running. The free triangle vanishes; the four-point
coefficient above follows from the 48 connected Wick pairings with the declared canonical normalization.
The full proof, exact controls and formal trace check accompany the updated register~\cite{ClosedSet}.}
\lines{26, 25}
\end{claim}

\textbf{Every compact simple group.} The lattice construction generalizes, but the determinant reduction and two-carrier
ordering are specific to $SU(2)$. Space remains $S^3$, and a dictionary starts the extension: $b_0\to11h^\vee/3$
($h^\vee$ the dual Coxeter number); the note $\tfrac34\to
C_{\min}(G)$, the smallest nontrivial Casimir; $\mathbb R^9/SO(3)\to\mathfrak g^3/\mathrm{Ad}(G)$; and the
Maurer--Cartan line runs through a long-root $\mathfrak{su}(2)$, where the minimal instanton lives~\cite{BPST75,BCGW77}.
I take $G$ simply connected; its other global forms share the local sector, which is all the gap and the local fields
see~\cite{AharonySeibergTachikawa13}.

\begin{claim}{PROVED (CITED)}{The lattice construction for compact groups}\label{b:everyGlat}
For the admissible compact-group lattice actions, reflection positivity, transfer-matrix construction and the
strong-coupling result hold with their stated hypotheses~\cite{OsterwalderSeiler78}.
\lines{1--3, 24}
\end{claim}

\begin{claim}{PROVED}{The electric note and the embedded doublet}\label{b:everyGmode}
For a compact connected simple Lie group on the Hamiltonian lattice (periodic side $\ge4$), the first electric level
is $2C_{\min}(G)g^2/a$, with the finite-volume qualification of Statement~\ref{b:kato}. In the two-dimensional continuum limit,
$\sigma_{\rm adj}/\sigma_{r_{\min}}\to h^\vee/C_{\min}$; and Helike's charge-2 doublet is an exact Hessian eigenvector
in the embedded long-root $\mathfrak{su}(2)$ block, with the $SU(2)$ frequency and normalization. This does not make
it the floor of the additional colour blocks.
\evidence{Peter--Weyl plus Gauss's law~\cite{KogutSusskind75}; the $j=1$
part of the adjoint under a long-root $\mathfrak{su}(2)$ is that $\mathfrak{su}(2)$~\cite{ClosedSet,ByEnough}.}
\lines{4--6, 10, 24}
\end{claim}

\textbf{The small room survives the change of colour group.} There was a precise gap in my dictionary:
the determinant estimate that confines the reduced potential was special to $SU(2)$. The brackets supply
its replacement. The global estimate I need is already in Austing and Wheater; scaling it closes this step.

\begin{claim}{PROVED}{The cubic--quartic bound and the reduced room}\label{b:cubicquartic}
For each fixed compact simple real Lie algebra with a fixed positive Ad-invariant inner product, put
\[
\mathcal C=\langle X_1,[X_2,X_3]\rangle,\qquad
Q=\sum_{i<j}\|[X_i,X_j]\|^2.
\]
There is a finite $c_{\mathfrak g}>0$, independent of the triple, such that
\[
|\mathcal C|\le c_{\mathfrak g}Q^{3/4},\qquad
V(q)\ge\frac\kappa2\|q\|^2-\frac{27S^4c_{\mathfrak g}^4}{32g^2}.
\]
Here $q=(X_1,X_2,X_3)$, $S=\sqrt t+2/\sqrt t$, $g>0$, and the original sources give
$\kappa=s[(t+2)-s(t+1)]/t>0$ for $t\ge1$, $0<s<(t+2)/(t+1)$.
The natural reduced Hamiltonian has compact resolvent, including on its colour-singlet charge sectors.
No group-uniform constant or lowest-carrier identification is asserted; this $c_{\mathfrak g}$ is not
the gap constant in H7 below.
\evidence{Deduction from the global bound of Austing--Wheater~\cite{AustingWheater03}, eq.~(2.2),
and the compact-form argument below; full convention audit in the cubic--quartic note accompanying~\cite{ClosedSet}.}
\lines{9 (extension), 24 (partial)}
\end{claim}

\emph{Proof.} In a faithful unitary representation $\rho$, let $A_k=i\rho(X_k)$ be Hermitian.
For some fixed $\eta>0$, the invariant inner product is
$\langle X,Y\rangle=-\eta\operatorname{Tr}(\rho(X)\rho(Y))$.
Thus $\mathcal C=-i\eta\operatorname{Tr}(A_1[A_2,A_3])$ and
$Q=\eta\sum_{i<j}\operatorname{Tr}([A_i,A_j]^\dagger[A_i,A_j])$.
The cited estimate translates to $|\mathcal C|\le A Q+B$ with fixed finite positive constants.
When $Q>0$, apply it to $rq$ with $r=Q^{-1/4}$:
\[
r^3|\mathcal C|\le A r^4Q+B=A+B,
\qquad |\mathcal C|\le(A+B)Q^{3/4}.
\]
When $Q=0$, all brackets and the cubic vanish. This is not an argument from compactness of $Q=1$;
that set is not compact. The global cited bound controls the nearly commuting directions.

Expanding the unchanged potential gives
$V=\tfrac12\sum_k(\omega_k^2-\mathrm{src}_k)\|X_k\|^2-gS\mathcal C+g^2Q/2$.
Its quadratic coefficients are at least $\kappa$. With $z=gQ^{1/4}/(Sc_{\mathfrak g})\ge0$,
\[
\frac{z^4}{2}-z^3+\frac{27}{32}
=\frac{(2z-3)^2(4z^2+4z+3)}{32}\ge0,
\]
which proves the stated lower bound. For the Friedrichs form on
$L^2(\mathfrak g^3,dq)$, bounded form norm controls the gradient and the quadratic spatial moment.
The latter controls the tails; Rellich compactness controls each bounded ball. The form embedding,
and hence the resolvent, is compact. Restricting to a closed reducing symmetry sector preserves compactness.
No boundary condition is added at a quotient singularity.\hfill$\square$

This removes one obstruction to changing groups. It does not order the new colour modes or remove the
continuum and large-room hypotheses. Discreteness was one requirement, not the whole extension.

\begin{claim}{CONJECTURE}{H7: the whole theory for every compact simple $G$}\label{b:h7}
For every compact simple $\mathfrak g$ there is $c_{\mathfrak g}>0$ such that, through the dictionary, the $\mathbb R^4$
theory exists with the OS axioms, has local fields governed by $b_0(G)$ at short distance, a unique vacuum, and
$\mathrm{spec}\,H\cap(0,c_{\mathfrak g}\Lambda_G)=\emptyset$.
\evidence{Clay$(G)$ rests on H2$(G)$, H3$(G)$, H4$'(G)$, H6$(G)$ only; the carriers are not needed, because the minimum
of continuous levels is continuous. Discreteness of the small room for general $G$ now follows from
Statement~\ref{b:cubicquartic}; the full H7 conclusion remains open. The scalar is lightest for $SU(3)$ too~\cite{MorningstarPeardon99}.}
\lines{24}
\end{claim}

The carriers are how I \emph{understand} the gap; they are not what Clay$(G)$ needs. What remains is to say in one
breath what the inventory is \emph{of}.
